% Part: second-order-logic
% Chapter: syntax-and-semantics
% Section: language-of-sol

\documentclass[../../../include/open-logic-section]{subfiles}

\begin{document}

\olfileid{sol}{syn}{lan}

\olsection{Basa Logika Orde Kapindho}

Kaya ing logika orde kapisan, ekspresi logika orde kapindho dibangun
saka tetembungan dhasar kang ngemot \emph{!!{variable}s},
\emph{!!{constant}s}, \emph{!!{predicate}s} lan kadhangkala
\emph{!!{function}s}. Saka unsur kasebut, bebarengan karo pangiket logis,
kuantor, lan tandha wacan kayata kurung lan koma, dibentuk
\emph{term} lan \emph{!!{formula}s}. Bedane yaiku saliyane variabel
kanggo objek, logika orde kapindho uga ngemot variabel kanggo relasi
lan fungsi, lan ngidini kuantifikasi marang variabel kasebut.
Dadi tandha logis ing logika orde kapindho yaiku tandha logis ing
logika orde kapisan, ditambah:

\begin{enumerate}
\item Himpunan !!{denumerable}s saka !!{variable}s relasi orde kapindho
  kanggo saben aritas~$n$: $\Obj V_0^n$, $\Obj V_1^n$, $\Obj V_2^n$, \dots
\item Himpunan !!{denumerable}s saka !!{variable}s fungsi orde kapindho:
  $\Obj u_0^n$, $\Obj u_1^n$, $\Obj u_2^n$, \dots
\end{enumerate}

Ing logika orde kapisan, !!{identity}~$\eq$ biasane kalebu.
Ing logika orde kapisan, tandha nonlogis saka basa~$\Lang{L}$
wigati supaya kita bisa ngandharake prakara kang migunani.
Mesthi ana !!{sentences} kang ora nggunakake tandha nonlogis,
nanging mung nganggo~$\eq$ angel ngandharake prakara kang migunani.
Ing logika orde kapindho, amarga variabel relasi lan fungsi kasedhiya
tanpa wates, kita bisa ngandharake kabeh kang bisa diandharake ing
basa orde kapisan sanajan tanpa nyedhiyakake tandha nonlogis khusus.

\end{document}
